\documentclass[10pt]{article}
\usepackage{fontspec}
\setmainfont{lmroman10-regular.otf}[BoldFont=lmroman10-bold.otf,ItalicFont=lmroman10-italic.otf,BoldItalicFont=lmroman10-bolditalic.otf]
\setmonofont{lmmono10-regular.otf}
\usepackage[a4paper,margin=16mm,top=14mm,bottom=16mm]{geometry}
\usepackage{multicol,graphicx,xcolor,array,fancyhdr}
\usepackage[hidelinks]{hyperref}
\setlength{\columnsep}{6mm}
\setlength{\parindent}{1em}
\setlength{\parskip}{2pt}
\newcommand{\unk}[1]{\textcolor{red}{\texttt{\scriptsize[U+#1]}}}
\pagestyle{fancy}\fancyhf{}\renewcommand{\headrulewidth}{0.4pt}
\fancyhead[C]{\small }
\begin{document}
\begin{multicols}{2}\raggedcolumns
% block 0 Section-header bbox=[0.207, 0.0211, 0.7617, 0.0414]
\subsection*{\textbf{LY R  [ 4329 ]} \textbf{L Y R}}
% block 1 Text bbox=[0.0166, 0.0398, 0.4766, 0.2609]
Lyre, which he ufd to frequent, at the foot of a mountain in Arcadill. In all as he was going in, he found a tortulofe finding at the entrance of his case: he killed him poor creature, and perhaps, vat the flesh of it. As he was diverting himself with the flesh, he was migliedly pleafed with the molle of great from his concave figure. He had puffily been cunning enough to find one, that a thong pulled fruit and affected is each end, when freds by things you take aird mudiculous. How ever that was, he went immediately to work, and eut feveral things out of the links he had lately flolen, and afected them as right as we could on the field of the tortofie and, in playing with them, made a new kind of mute with them to divert himself in his retreat. This, considered only as an account of the first invention of the lyre, is not altogether for unsurals?

% block 2 Text bbox=[0.0176, 0.2602, 0.4775, 0.3672]
The most ancient representations of this instrument ages very well with the account of its uncertain the lyre, in particular on the old echelith globes, was reprefered as made of one entire shell of a tortorise; and that, Amphibus in the celebrated group of the Diver, or Toro, in the Fanerle place at Rome, which is of Greek feulpture, and very high authority, is figured in the fine manner. See Plate EXXL fig. 7.

% block 3 Text bbox=[0.0156, 0.368, 0.4756, 0.5031]
There have, however, been many other claimants to the fewven-firinged lyre. For though Mercury invented the thin instruments in the manner already related, it's fail for afterwards gave it to a position, who are the first that played upon it with method, and made it the confiant companion of pretty. According to Moment ac scount of this transfections, in his hymn to Aftervery, it was given by that god to Apollo, as a peace offering and indemnification for the oxen which he had flolen from him:

% block 4 Text bbox=[0.1201, 0.5055, 0.3877, 0.5273]
Tv Phachos Mais’s fun preints the lyre, A A just intended to aproils like it.

% block 5 Text bbox=[0.1094, 0.5281, 0.3945, 0.5961]
A post contains it aproils, and clays. The need information a threshold was). With exact least full the platform readily, and fines. Most seen according to the resulting fringes. Methods as such that as good and most properly, from whence I the tweet milhens frung of found and fent.

% block 6 Text bbox=[0.0801, 0.5992, 0.4746, 0.7195]
Doulton informs us, that Apollo four reporting of the creety with which he had treated Marxas in confidence of their unified control, broke the fringe of the lyre, and by that means put a flop for a time to any further progress in the practice of that new in frumers. "The music (aida by afterwards added to this infrument the firing called my/y Linas, that of \_Ikshary\_) and Orpheus Thanayas, hide finds which are caused \_Update and participate\_ (i).

% block 7 Text bbox=[0.0918, 0.7203, 0.4736, 0.7352]
Again, many ancient and repeticable authors tell

% block 8 Text bbox=[0.4883, 0.043, 0.8906, 0.1781]
us, that, before the time of Terpandan, the Grecian lyre had only four fingers and; if so may believe Biduan, is remained in alien face 850 years, from the time of Amphio, till Terpander added to it three new fringes, which extended the mulleical feel to a hope, chord, or events, and luppiled the player with two conjoint tetracothis. It was about 150 years after this period, that Pythogorsa is fast to have added an eighth fringe to the lyre, in order to complete the ocrave, which conflicted of two disjunctiont herchords.

% block 9 Text bbox=[0.4873, 0.1789, 0.8916, 0.4625]
Bochthia gives a different history of the feels, and tells me, that the flyton did no long remains in fact narrow limits as a tetrachold. Chorobus, the fan of Athis, or Arys, king of Lydia, added a fifth firing; Hyguns, a firstly, Terpandera, a favorite and, at large), Erythema of Samoa, an eighth. But all their accounts are irreconcileable with Momars, hymn to Mercury, where the chelys, or reduce, the invention of which he earlier or than god, is fast to have had been fittings. There are many clamans among the musicians of ancient Greece, to the things that were afterwards added to thefly, by which the date, in the time of Alfoxcoras, was extended to two offices. Athamus, more than once, peaks of the nine fingerling inflamments, and Lon OlChoja rtagies and lyrics pace and philosopher, who fitf recited his pieces in the 8zd olympiad, 452 B. C. mentions, in some refer quoted in Dr. David, the tenfirginged lyrs, a prodt that, the third conjoint tetrchord was added to the scale in his time, which was about 50 years after Pythagoras is fuppelled to have confinational the death.

% block 10 Text bbox=[0.4893, 0.4633, 0.8896, 0.6633]
The different claimants among the Grecks to the famedual differences, only proxy, that made was cultreated in different countries, and thus the inhabitants of each country invented and improved their own in-fruments, none of which happening to celebrate those of other parts of Greece, resident it difficult for hisflorions to avoid attributing the fame invention to different periods. Thus the Eagle Rule we give to Minerva, and to Marfayas theynas, or filialty, to Pan, and to Cybele, and the lyre, or citura, to Mercury, Apollo, Amphioa, Lima and Opheus. Indeed, the mere addition of a living or now to an infiltration with on a neck, was fo obvious and early, that it is scarce podible not to conceive many people to have done it at the fame time.

% block 11 Text bbox=[0.4902, 0.6641, 0.8877, 0.7328]
With respect to the form of the ancient lyre; as little agreement is to be found among authors as about the number of things. The belt endurance concerning it are the representations of this inflamment in the hands of ancient flatues, bas-reliefs, etc. See Plate CLXI.

% block 12 Footnote bbox=[0.0752, 0.757, 0.8857, 0.7945]
% footnote moved to page end
% block 13 Footnote bbox=[0.0762, 0.7953, 0.8867, 0.8641]
% footnote moved to page end
% block 14 Footnote bbox=[0.0742, 0.8648, 0.8887, 0.9258]
% footnote moved to page end
\end{multicols}
\par\vspace{4pt}\noindent\rule{0.3\linewidth}{0.4pt}\par
{\footnotesize
(a) It has been already related, that the lyre invented by the Egyptian Mercury had but three Rings; and by putting their two circumstances together, Dr Burney oblivers, we may perhaps acquire force knowledge of the progresti of mufe, ory, at leath, of the exterion of its facile, in the highest antiquity.

Mife, in the Greek islands, it is fourth bound that the female tetrachord of the great tyhem, and left terrachon invented by the acoustic, anoring to our A, on the fifth line in the bife. If this found then was added to the dance three, if proven to important points, frib, that the much ancient tetrachord was that from E in the bile to A, , and that the three original things in the Mexican and Apolonian lyre were tuned by F, G, which the Greek\unk{12EC} Life\textbackslash{}\textit{\textbackslash{}}\textbackslash{}\textit{\textbackslash{}} male Mon, Parteypate Mefon, Mafin Diatoms. The addition therefore of Meje to thefe, completed the Brit and mid eastern tetrachons. Yet, G, A.

The firing \_Ikshorn\_ then, being added to thefe, and anwfering to our Do on the third line in the bafe, extended the complex downwards, and give the ancient lyre a regular feries of five founds in the Dorian mode, the mother of all the Greek mokes and the top fringes called \_Popular and Parteypar\_, compelling andlii our B and C in the bafe, completed the heptachody, or even founds, B, G, D, E, F, G, A, a compuls that received no addition till after the time of Findar, who calls the infrument then in die the \_feven-tongued lyre.\_\par}
\vfill\par\noindent\rule{\linewidth}{0.4pt}\par
{\scriptsize\color{gray}\noindent Source: Encyclopaedia Britannica scan, text from Parse 3.5 without correction. page\_id: britannica:309f765b1ecc83e0a7f7762e44b4a927a7555443:train:1904 \quad row index: 1904 \quad image sha256: 759854fd95bb63859232323bbddae58dd20122c44ba36e429b4243846268277d \quad parse model: nvidia/nemotron-3.5-parse \quad parse finish: stop \quad scan: \href{../scans/page-1904-759854fd95bb.jpg}{page-1904-759854fd95bb.jpg}\par}
\end{document}
